\documentclass[11pt,a4paper]{article}
\usepackage[utf8]{inputenc}
\usepackage[T1]{fontenc}
\usepackage[french]{babel}
\usepackage{amsmath, amssymb, amsthm, mathrsfs}
\usepackage{hyperref}
\usepackage{geometry}
\geometry{margin=1in}
\usepackage{fancyhdr}

\pagestyle{fancy}
\fancyhf{}
\fancyfoot[C]{\thepage}
\fancyfoot[R]{\small Charles EDOU NZE, chercheur indépendant}

\newtheorem{theorem}{Théorème}[section]
\newtheorem{lemma}[theorem]{Lemme}
\newtheorem{definition}[theorem]{Définition}
\newtheorem{corollary}[theorem]{Corollaire}
\newtheorem{conjecture}[theorem]{Conjecture}
\newtheorem{remark}[theorem]{Remarque}

\title{Approche Rigoureuse de la Conjecture d'Erdős-Straus}
\author{Charles EDOU NZE}
\date{\today}

\begin{document}
\maketitle

\begin{abstract}
La conjecture d'Erdős-Straus postule que pour tout entier $n \geq 2$, il existe des entiers strictement positifs $x, y, z$ tels que $4/n = 1/x + 1/y + 1/z$. Nous fournissons des définitions axiomatiques de l'équation diophantienne, examinons les développements récents, isolons des lemmes basés sur l'arithmétique modulaire, et esquissons une architecture formelle pour une autoformalisation potentielle dans Lean 4.
\let\thefootnote\relax\footnotetext{Charles EDOU NZE, chercheur indépendant}
\end{abstract}

\section{Définitions Axiomatiques et Énoncé du Problème}
Soit $\mathbb{N}^*$ l'ensemble des entiers strictement positifs, $\mathbb{N}^* = \{1, 2, 3, \ldots\}$.
\begin{definition}[Représentation d'Erdős-Straus]
Pour un entier $n \in \mathbb{N}_{\geq 2}$, une \textit{représentation d'Erdős-Straus} de $n$ est un triplet $(x, y, z) \in (\mathbb{N}^*)^3$ satisfaisant l'équation diophantienne :
\begin{equation} \label{eq:es}
\frac{4}{n} = \frac{1}{x} + \frac{1}{y} + \frac{1}{z}
\end{equation}
\end{definition}

\begin{conjecture}[Conjecture d'Erdős-Straus]
Pour tout entier $n \ge 2$, il existe au moins une représentation d'Erdős-Straus $(x, y, z) \in (\mathbb{N}^*)^3$.
\end{conjecture}

Il est une réduction fondamentale bien établie que si la conjecture est vraie pour tous les nombres premiers $p \ge 3$, elle est vraie pour tous les entiers $n \ge 2$. Ainsi, nous restreignons notre domaine à $n = p \in \mathbb{P}_{\ge 3}$.

\section{Recherche de Littérature Contextuelle}
Une ligne importante de recherche récente sur la conjecture implique le paramétrage des solutions basé sur des classes de congruence. Notamment, l'article «~Congruence Classes of Supporting the Erdös-Straus Conjecture I: Tame Solutions~» explore la nature algébrique des solutions lorsque $n = 24m + 1$. Une solution $(n_1, n_2, n_3)$ avec $n_1 \le n_2, n_3$ où $n_1 = 6m + k$ est qualifiée de \textit{solution docile} (tame) si $n_2$ et $n_3$ divisent $(6m+k)(24m+1)$. La littérature établit que parmi les nombres premiers de la forme $24m+1$ avec $m \le 30000$, les nombres premiers sauvages (ceux dépourvus de solutions dociles) sont exceptionnellement rares. Cette catégorisation structurelle motive notre recours à l'arithmétique modulaire pour partitionner l'espace des solutions premières.

\section{Stratégie de Preuve et Isolation des Lemmes}
Pour aborder la conjecture générale sur les nombres premiers, nous décomposons les entiers modulo $24$. Les nombres premiers se répartissent en classes modulo $24$.
Nous définissons d'abord des identités polynomiales qui couvrent diverses classes de congruence.

\begin{lemma}[Identité Polynomiale Modulo la Classe] \label{lem:poly}
Soit $p$ un nombre premier. Si $p \equiv 3 \pmod 4$ ou $p \equiv 2 \pmod 3$, alors $p$ admet une représentation d'Erdős-Straus.
\end{lemma}

\begin{proof}
Nous analysons les cas par l'arithmétique modulaire. Nous cherchons à trouver $x, y, z$ sous forme de polynômes en $p$. Une transformation standard de l'équation \eqref{eq:es} donne $4xyz = p(xy + yz + zx)$.
Posons $x = p \cdot a$, $y = p \cdot b$, $z = c$.
L'équation devient :
\begin{equation}
\frac{4}{p} = \frac{1}{pa} + \frac{1}{pb} + \frac{1}{c} \implies \frac{4c - p}{pc} = \frac{a+b}{pab}
\end{equation}
Si nous choisissons $c = \lceil p/4 \rceil$, alors $4c - p \in \{1, 2, 3\}$.
Nous pouvons analyser les classes de résidus :

\textbf{Cas 1 :} $p \equiv 3 \pmod 4$. Alors $p = 4k + 3$ pour un certain $k \in \mathbb{N}$.
Posons $c = k + 1$. Alors $4c - p = 4(k+1) - (4k+3) = 1$.
L'équation se simplifie en $\frac{1}{p(k+1)} = \frac{a+b}{pab}$.
Nous pouvons satisfaire cela en choisissant $a = k+1+1 = k+2$ et $b = (k+1)(k+2)$.
Vérifions explicitement :
\begin{align*}
\frac{1}{p(k+2)} + \frac{1}{p(k+1)(k+2)} &= \frac{k+1 + 1}{p(k+1)(k+2)} \\
&= \frac{k+2}{p(k+1)(k+2)} \\
&= \frac{1}{p(k+1)} = \frac{1}{pc}
\end{align*}
Puisque $c = (p+1)/4$, nous avons :
\begin{align*}
\frac{1}{pa} + \frac{1}{pb} + \frac{1}{c} &= \frac{1}{pc} + \frac{1}{c} \\
&= \frac{1+p}{pc} \\
&= \frac{1+p}{p(p+1)/4} = \frac{4(p+1)}{p(p+1)} = \frac{4}{p}
\end{align*}
Chaque étape et changement d'indice est strictement positif et borné. Ainsi, pour $p \equiv 3 \pmod 4$, une représentation valide est explicitement générée par le triplet $(p(k+2), p(k+1)(k+2), k+1)$, où $k = (p-3)/4$.

\textbf{Cas 2 :} $p \equiv 2 \pmod 3$. Alors $p = 3k + 2$ pour un certain $k \in \mathbb{N}$.
Revenons à l'équation originale $\frac{4}{p} = \frac{1}{x} + \frac{1}{y} + \frac{1}{z}$. Posons $x = p$, ce qui est permis puisque $p \in \mathbb{N}^*$.
Alors $\frac{4}{p} = \frac{1}{p} + \frac{1}{y} + \frac{1}{z}$, ce qui signifie que nous devons résoudre $\frac{3}{p} = \frac{1}{y} + \frac{1}{z}$.
Posons $y = p \cdot m$ et $z = p \cdot m \cdot n$. Alors $\frac{3}{p} = \frac{1}{pm} + \frac{1}{pmn}$, ce qui implique $3mn = n + 1$.
Au lieu de cette substitution, multiplions $\frac{3}{p} = \frac{1}{y} + \frac{1}{z}$ par $pyz$. Nous obtenons $3yz = p(y+z)$.
Posons $y = c$ et $z = pd$. Alors $3cpd = p(c+pd) \implies 3cd = c + pd$.
Isolons $c$ : $c(3d - 1) = pd$. Ainsi $c = \frac{pd}{3d-1}$.
Pour s'assurer que $c$ est un entier, nous avons besoin que $3d - 1$ divise $pd$. Puisque $p$ est premier, si $3d - 1$ divise $p$, alors $3d - 1 \in \{1, p\}$.
Si $3d - 1 = 1$, alors $3d = 2$, ce qui n'a pas de solution entière.
Si $3d - 1 = p$, alors $3d = p + 1$. Puisque $p = 3k + 2$, nous avons $p + 1 = 3k + 3 = 3(k+1)$.
Ainsi, nous pouvons choisir $d = k + 1$.
Alors $c = \frac{p(k+1)}{p} = k + 1$.
Vérifions explicitement avec ces choix. Nous fixons $x = p$, $y = k + 1$, et $z = p(k+1)$.
\begin{align*}
\frac{1}{x} + \frac{1}{y} + \frac{1}{z} &= \frac{1}{p} + \frac{1}{k+1} + \frac{1}{p(k+1)} \\
&= \frac{k+1}{p(k+1)} + \frac{p}{p(k+1)} + \frac{1}{p(k+1)} \\
&= \frac{p + k + 2}{p(k+1)}
\end{align*}
Substituons $k = (p-2)/3$ :
\begin{align*}
\frac{p + (p-2)/3 + 2}{p((p-2)/3 + 1)} &= \frac{(3p + p - 2 + 6)/3}{p(p - 2 + 3)/3} \\
&= \frac{4p + 4}{p(p+1)} \\
&= \frac{4(p+1)}{p(p+1)} = \frac{4}{p}
\end{align*}
Puisque $p \ge 2$, $k \ge 0$, et donc $x, y, z$ sont des entiers strictement positifs. Par conséquent, pour $p \equiv 2 \pmod 3$, une représentation valide est explicitement générée par le triplet $(p, (p+1)/3, p(p+1)/3)$.

Ceci conclut la preuve constructive pour ces deux classes de congruence spécifiques.
\end{proof}

\section{Architecture pour l'Autoformalisation}
Pour faciliter la vérification symbolique future via un démonstrateur de théorèmes agentique tel que Lean 4, nous définissons les types formels et les hypothèses nécessaires.

\begin{verbatim}
import Mathlib.Data.Real.Basic
import Mathlib.Data.Nat.Prime

-- Definition of the Erdős-Straus representation
def is_es_representation (n x y z : \mathbb{N}) : Prop :=
  x > 0 \land  y > 0 \land  z > 0 \land  (4 : \mathbb{Q}) / n = 1 / x + 1 / y + 1 / z

-- Reduction Lemma signature
theorem es_reduction_to_primes (h : \forall  p : \mathbb{N}, p.Prime \to
  \exists  x y z, is_es_representation p x y z) :
  \forall  n : \mathbb{N}, n \ge  2 \to  \exists  x y z, is_es_representation n x y z := by
  sorry

-- Polynomial Modulo Lemma
theorem es_mod_4_case_3 (k : \mathbb{N}) :
  let p := 4 * k + 3
  is_es_representation p (p * (k + 2)) (p * (k + 1) * (k + 2)) (k + 1) := by
  sorry
\end{verbatim}

\end{document}
